\chapter{Talking to Language Models}
\label{ch:llm}

\standfirst{A model behind an API, asked from a workspace: one prompt and
its answer, the answer as it arrives, a conversation that remembers, tools
the model runs through you, and the vector a text becomes. Four services,
one way of asking.}

\section{What you need}

A key, and a build that speaks TLS.

The key goes in the environment before \ct{st80} starts, because that is
where a secret belongs and the one place this system reads one from:

\begin{sh}
export ANTHROPIC_API_KEY=sk-ant-...
./st80
\end{sh}

\begin{st}
Smalltalk environmentAt: 'ANTHROPIC_API_KEY'
\end{st}

answers it, or \ct{nil} when it was never set.

TLS is OpenSSL, found when the system was built and optional the way ODBC
is (Chapter~\ref{ch:building}): \ct{make deps} says whether the machine has
it, \ct{Socket isTlsAvailable} says whether this image does, and an
\ct{https} URL on a build without it is a \ct{NetError} saying so by name.
Every cloud service here is https only. A local Ollama is plain http and
needs neither the key nor the library.

\section{The first call}

\begin{st}
Smalltalk at: #Key put: (Smalltalk environmentAt: 'ANTHROPIC_API_KEY').
Smalltalk at: #Claude put: (Anthropic apiKey: Key model: 'claude-opus-5').
Claude send: 'What is the capital of France?'
\end{st}

The answer is a String, the whole of it, once the model has finished.
\ct{send:} is the same message on every service; only the first line
differs:

\begin{st}
Smalltalk at: #Gpt put: (OpenAI apiKey: Key model: 'gpt-4o').
Smalltalk at: #Router put: (OpenRouter apiKey: Key model: 'anthropic/claude-sonnet-4').
Smalltalk at: #Llama put: (Ollama new model: 'llama3.2'; yourself).
\end{st}

What the model is told beside the prompt is set beforehand, and only what
was set is sent---a model runs at its own defaults until you say
otherwise:

\begin{st}
Claude systemPrompt: 'Answer in one sentence.'; temperature: 0.2; maxTokens: 400.
Claude timeoutSeconds: 120.                     "five minutes unless told"
\end{st}

\ct{temperature:} and \ct{topP:} forget each other, because Anthropic
refuses a request naming both. \ct{maxTokens:} is sent when set, and
always to Anthropic, which requires it: 8192 unless told.

Three things go wrong, three ways. A reply that is not 2xx is an
\ct{LLMError} carrying the status and the body, which is the service's own
sentence about what was wrong---a bad key is a 401 that says so. A service
that cannot be reached is a \ct{NetError}. One that takes longer than the
timeout is a \ct{TimedOut}.

\begin{st}
[Claude send: 'hi'] on: LLMError do: [:e | e status printString, ' ', e body]
\end{st}

And what the last call left behind is kept: \ct{lastResponse} is the whole
JSON of a reply that was not streamed, for the token counts and whatever
else this does not surface; \ct{lastStopReason} is why the model stopped,
in the service's own word---\ct{end_turn}, \ct{stop}, or \ct{refusal},
which is not an error but an empty answer with the reason in
\ct{lastResponse}.

\section{The answer as it comes}

\begin{st}
Claude send: 'Tell me a story' do: [:piece | Transcript show: piece]
\end{st}

Each piece of text goes to the block as the service sends it---a few words
at a time, as server-sent events from Anthropic and OpenAI and as one JSON
document a line from Ollama---and the whole is answered at the end. The
wait is the same wait as before, and parks the process and not the worker,
so a server answering many requests can have many of these in flight on a
few threads. The timeout bounds each piece rather than the whole: a model
that keeps talking is allowed to.

\section{A conversation}

A model remembers nothing between calls; a conversation does, and sends
the whole of what was said with every question:

\begin{st}
Smalltalk at: #Talk put: Claude conversation.
Talk ask: 'My name is Ada.'.
Talk ask: 'What is my name?'.
Talk ask: 'Say it slowly' do: [:piece | Transcript show: piece].
Talk messages size.                             "6: three questions, three answers"
Talk clear
\end{st}

\ct{messages} is the history in the service's own shape---Anthropic's
content blocks, OpenAI's tool calls---because a tool exchange has to be
echoed back exactly as the service wrote it, and because the shape is the
model's business. A conversation belongs to one model; to talk to another,
start another. \ct{send:} is a conversation of one question.

\section{Tools}

A tool is something the model may ask you to run: a name, the description
it reads to decide when, a JSON schema for the arguments it should send,
and a block that does it. The block gets the arguments as a
\ct{JSONObject} and answers what the model should read---a String as it
is, JSON as JSON, anything else printed.

\begin{st}
Smalltalk at: #Now put: [:args | Time now printString].
Talk addTool: (LLMTool name: 'time' description: 'the time now' do: Now).
Talk ask: 'What time is it?'
\end{st}

With arguments, the schema says what they are:

\begin{st}
Smalltalk at: #Schema put: '{"type":"object","properties":{"n":{"type":"number"}}}' asJson.
Smalltalk at: #Sq put: [:args | (args numberAt: 'n') squared printString].
Talk addTool: (LLMTool name: 'square' description: 'n squared' parameters: Schema do: Sq).
Talk ask: 'Use the square tool on 12 and answer with the number alone.'
\end{st}

What happens underneath is a loop. The question is posted with the tools;
when the reply asks for one, the reply's own message is added to the
conversation, the block runs, its answer is added as the tool's result in
the service's shape, and the model is asked again---until it answers with
text, which is what \ct{ask:} answers. Twenty rounds is the limit, after
which it is an \ct{LLMError} and not a conversation. A model that asks for
a tool nobody gave it is told \ct{no such tool} as the tool's result and
asked again, since a model that invents a tool reads the answer and gets
over it.

A tool reads its arguments defensively. The schema says \ct{n} is a
number, and a large model sends one; a small local model sends \ct{'12'}
as often as not, so the robust block takes \ct{args at: 'n'} and converts
a String rather than asking \ct{numberAt:} and raising.

Tools belong to the conversation, not the model: the same model answers
plainly in one conversation and runs tools in another. And a conversation
with tools answers whole rather than as it comes, even to \ct{ask:do:}: a
reply that turns out to be a tool call has no text to stream, and the
request cannot know in advance, so the block gets the text once, at the
end.

\section{A picture}

\begin{st}
Claude send: 'What is in this chart?' image: 'chart.png'
\end{st}

The file is read as bytes and sent inline, base64, with its type by the
extension---PNG, GIF, WebP, and JPEG for the rest---in each service's own
shape: a content block for Anthropic, a data URL for OpenAI and OpenRouter,
an \ct{images} list for Ollama. \ct{ask:image:} does the same in a
conversation, and OpenAI's \ct{imageDetail:}---\ct{low}, \ct{high},
\ct{auto}---says how closely to look.

\section{Embeddings, and where to keep them}

\begin{st}
Smalltalk at: #Embedder put: (OpenAI apiKey: Key model: 'text-embedding-3-small').
(Embedder embed: 'Cats sleep sixteen hours a day.') size          "1536"
\end{st}

\ct{embed:} answers an Array of Floats---Floats, and not the exact
Fractions the JSON reader answers for decimals, which are right for money
and wrong for fifteen hundred coordinates. OpenAI and OpenRouter have an
embeddings endpoint; Ollama has one for its embedding models; Anthropic has
none and says so.

Where they go is \ct{Qdrant}, a vector database, asked over plain HTTP on
this machine:

\begin{st}
Smalltalk at: #Store put: (Qdrant on: 'notes').
Store createCollection: 'notes' size: 1536.
Smalltalk at: #Text put: 'Cats sleep sixteen hours a day.'.
Store store: (Embedder embed: Text) text: Text document: 'notes' sequence: 1 payload: nil.
(Store search: (Embedder embed: 'how long do cats sleep?') limit: 3) asJsonString
\end{st}

A record is a vector, an id this makes when none is given, and a payload
carrying the text, the document it came from and its place in it, plus
whatever \ct{payload:} adds; a search answers the nearest records as the
server's own JSON, each with its id, its score and its payload. That is
the whole of a retrieval-augmented answer: embed the question, search, put
the texts found in front of the model, ask.

\section{The four services}

\begin{center}
\small
\begin{tabular}{@{}llp{6.2cm}@{}}
\toprule
class & the key, and where it talks & what is its own\\
\midrule
\ctp{Anthropic} & \ctp{ANTHROPIC\_API\_KEY}, \ctp{api.anthropic.com} & \ctp{maxTokens:} always sent; \ctp{effort:} how hard it thinks; \ctp{fallbacks:} models to try when it declines; \ctp{version:}\\
\ctp{OpenAI} & \ctp{OPENAI\_API\_KEY}, \ctp{api.openai.com} & \ctp{reasoningEffort:} for the models that reason; \ctp{imageDetail:}; \ctp{embed:}\\
\ctp{OpenRouter} & \ctp{OPENROUTER\_API\_KEY}, \ctp{openrouter.ai} & many services' models as \ctp{provider/model}; \ctp{fallbackModels:} tried in order; \ctp{siteUrl:} and \ctp{siteName:} for the gateway's credit\\
\ctp{Ollama} & no key; \ctp{localhost:11434} & \ctp{isUp}, never an error; \ctp{models}; \ctp{url:} for another machine; \ctp{embed:} with an embedding model; \ctp{Ollama toHtml:} for a screen\\
\bottomrule
\end{tabular}
\end{center}

\ct{url:} on any of them points it at a proxy, a gateway or a test. Every
one of them is asked the same way; the class knows only its wire: the
endpoint and the headers, the shape of a message and a tool, where the text
and the tool calls are in a reply, and how a streamed reply is cut up.
\ct{doc/LLM.md} lays the four wires side by side.

\section{TLS, and what it refuses}

An \ct{https} URL connects and then makes the socket the client end of a
TLS connection. The certificate is checked against the system's own store
and its name against the host you asked for, as a browser checks them, and
neither check can be turned off---a client that can be asked to trust
anything will one day be asked to. Anything wrong is a \ct{NetError} that
says what:

\begin{st}
(HttpClient get: 'https://api.anthropic.com/v1/models') status         "401: a reply"
[HttpClient get: 'https://expired.badssl.com/'] on: NetError do: [:e | e messageText]
\end{st}

\begin{plain}
'TLS handshake: certificate verify failed: certificate has expired'
\end{plain}

and likewise \emph{hostname mismatch}, \emph{self-signed certificate},
\emph{IP address mismatch} for an address that is not in the certificate,
and \emph{wrong version number} from a peer that does not speak TLS at
all. A private authority goes in the system's store, or in
\ct{SSL_CERT_FILE} in the environment, which OpenSSL reads.

\begin{stwarn}[The server side is not here]
This system's servers speak plain HTTP. A server on the internet sits
behind a reverse proxy---nginx, Caddy---that terminates TLS and passes the
request on, which is the arrangement Kiss has with Tomcat and the one
Chapter~\ref{ch:serving} assumes. TLS here is the client's, for reaching
services that are https only.
\end{stwarn}

Underneath, the socket stays non-blocking and the arm-wait-retry contract
of Chapter~\ref{ch:serving} stays the whole contract; what TLS adds is that
a read may have to write first and a write may have to read first, and the
primitive says which way to wait. \ct{doc/NETWORK.md} has it.

\section{Testing it}

Thirty-three tests in \ct{lib/LLM-Tests} reach no service. They run each
class against a listener of the test's own that answers canned replies
and keeps every request it was sent, parsed---so what is checked is the
wire: the header each service wants, the body a request carries and does
not carry, a streamed reply cut into pieces with the chunk boundary falling
inside an event, a tool call and its answer going back in the second
request, an embedding, and what a refusal, a rate limit and an error
mid-stream become. Two profiles reach the real thing and are run on
purpose: \ct{internet-live} fetches real certificates---Anthropic's without
a key, for a 401 over a good one, and three from badssl.com that must be
refused---and \ct{llm-live} asks the four services with keys from the
environment, failing by name for a key that is not there.

\begin{stnote}[Found while building it]
\ct{#(true false)} is two Symbols in this dialect, not two Booleans, and a
test that meant the Booleans said so before it passed. A method may
reference sixty-three literals and no more, in the 1983 compiler; one test
was split. The JSON reader's exact decimals---\ct{0.25} is \ct{1/4}---are
the wrong thing for a vector, and \ct{embed:} converts. An event stream's
lines cross chunk boundaries, because a server chunks by its buffer and
not by the line; the body reader reads through them, and the test puts a
boundary inside an event to prove it. And a peer that accepts a connection
and then says nothing holds a TLS handshake for ever, which a client's read
timeout did not cover until it did.
\end{stnote}

\section{Where the rest of it is}

\ct{doc/LLM.md} is the design: what \ct{LLM} knows once and what each
subclass answers, the four wires in one table, the tool loop, and what is
deliberately not here---OpenAI's Responses endpoint, Kiss's retries,
server-side TLS. \ct{lib/LLM/PROVENANCE.md} records what crossed from
Kiss's \ct{org.kissweb.llm} and what changed. \ct{doc/HTTP-CLIENT.md} is
the client underneath: https, and a reply read as it comes.
\ct{doc/NETWORK.md} is TLS in the socket layer. The tests are
\ct{lib/LLM-Tests}, \ct{lib/HTTP-Client-Tests} and \ct{lib/Network-Tests};
the live ones \ct{lib/LLM-Live-Tests} and \ct{lib/HTTP-Client-Live-Tests}.
